% !TEX root = ../main.tex
\chapter{Review: judge what it found}\label{ch:review}
Review is where a person gives a \term{verdict}{verdict} on what the machine proposed. A
detection run, a seed search, a window set or a Library family each produce
\term{candidate}{candidates}; Review shows them to you one at a time, you press one key per
candidate, and your answer goes straight into the database.

Everything in Review is organised by \term{review-queue}{queue}. A queue is a standing
question about one source, such as ``every detection of run 32'' or ``every span taken for
Review in Explore''. It holds no list of its own: each time you open it, it asks the source
again, so a verdict given anywhere else in the app shows up here at once.

There are two pages. \apppage{review.inspector} judges one candidate at a time and is where
almost all reviewing happens. \apppage{review.cluster} judges a group of similar candidates
in one stroke. A third view, the blind test, appears in place of the inspector for one kind
of queue (\cref{sec:review-blind-test}).

\section{The five verdicts}\label{sec:review-verdicts}
The verdict vocabulary is fixed at five words. The app, the database and every count in
the manual use exactly these, and no page adds a sixth.

\begin{description}
  \item[seed (key S)] ``This is a clear, typical example of something I want to find more
    of.'' Use it sparingly, for the best example of a kind of event. A seed does more than
    record an opinion: it \emph{promotes} the candidate into the Library as a new
    \term{motif}{motif} entry, which Discovery can then search the recordings for.
  \item[interesting (I)] ``This is a real event worth keeping.'' The everyday yes.
  \item[not interesting (N)] ``There is nothing here worth keeping.'' The everyday no.
  \item[artifact (A)] ``This is not the fungus: it is the equipment.'' A loose electrode, a
    step in the baseline, a wire picking up its neighbour. Artifacts are kept out of
    training data.
  \item[unsure (U)] ``I cannot tell.'' Use it honestly rather than guessing; an unsure
    verdict still counts as judged, so the item will not come back.
\end{description}

\ui{Space} skips a candidate and writes nothing. It is not a verdict.

\begin{plainterms}[Why your verdict never edits a detection]
The app keeps two kinds of record in separate tables: \term{detection}{detections}, which only
the machine writes, and \term{annotation}{annotations}, which only people write. Think of a
lab notebook in which the instrument's printout is pasted in and never written on; your
comments go on the facing page. When you judge a detection, your verdict is stored as a
separate row (an \emph{adjudication}) that points at the detection. The detection itself
is untouched, so the run that produced it can always be re-run and checked, and a later
reader can always tell what the machine said from what a person said. Each queue states,
in its \ui{writes} chip, which table its verdicts go to.
\end{plainterms}

\pagesection{review.inspector}{Queue}{One candidate at a time}{live}
The inspector shows one candidate with everything you need to judge it: the signal around
it, its shape, how the machine scored it, and the five verdict buttons. You judge, the page
moves on to the next unjudged candidate, and you judge again.

\pagefield{When you'd come here}
Whenever a queue has candidates waiting: after a Discovery run sends its detections, after
you take a span for Review in Explore, after a training run sends windows, or when a Library
family sends its unjudged members. \unverified{The \ui{need you} count in the header shows how many are
waiting across all queues.} Opening Review on its own lands on the last queue you used, or
the first queue there is.

\screenshot{pages/review/review.inspector--1-candidate (deep link, live detection)}{A
detection from a live queue. The toolbar names the queue and its write target; the
context card shows the candidate shaded in its surroundings; the verdict row sits below the
shape.}

\pagefield{Doing things here}
\subsection{Judge a queue}\label{sec:review-judge}
\begin{enumerate}
  \item Choose a queue with the \ui{queue} chip at the top left. Each queue shows how many
    candidates are left.
  \item Read the candidate. The meta line under its id gives the recording, channel, hours
    and length. \emph{Candidate in context} shows it shaded with $\pm$30, 120 or 300\,s
    around it (\ui{padding}). \emph{Shape} draws it alone in millivolts, at its real
    resolution: a short event at 1\,Hz looks like a staircase because it is one. When the
    recording was made by down-sampling a faster one, \ui{source resolution} draws the same
    stretch from the faster original.
  \item Check the neighbours if you suspect an artifact: \ui{Other channels} shows the same
    stretch on every other channel of the recording.
  \item Give a verdict with a key (\ui{S}, \ui{I}, \ui{N}, \ui{A}, \ui{U}) or a button. The
    verdict is written to the database first; only then does the page mark it and, with
    \ui{auto-advance} on, move to the next unjudged candidate.
  \item Made a mistake? \ui{Ctrl Z} undoes the last write in this queue, on the server: the
    database restores whatever was there before, including an earlier verdict if you had
    changed your mind. \ui{Ctrl Shift Z} redoes it.
\end{enumerate}

\subsection{Classes, tags and notes}\label{sec:review-annotate}
The \emph{Annotate} card is optional. The number keys pick a class: \ui{1} spike-train,
\ui{2} burst, \ui{3} slow-drift, \ui{4} plateau, \ui{9} electrode artifact. A class also gives
a verdict: an informative class means \emph{interesting}, and electrode artifact means
\emph{artifact}; that verdict is what is written. Tags must come from the project's tag
vocabulary (Settings › Vocabulary); a term outside it is refused, together with the verdict.
Tags and the note travel with the next verdict you give.

\subsection{Promote a seed}\label{sec:review-promote}
Pressing \ui{S} writes the seed verdict and, in the same step, creates a Library entry for
the candidate's shape, or joins an existing entry if the same shape is already held.
Auto-advance pauses and the \emph{Promoted to Library} panel appears. It names the entry,
offers \ui{Open in Library} and \ui{Seed search in Discovery}, and asks which family the new
exemplar belongs to. \ui{Confirm and next} (Enter) moves on; \ui{Undo promotion}
(\ui{Ctrl Z}) removes both the verdict and the Library rows it created.

\subsection{Read the evidence}\label{sec:review-evidence}
The right-hand rail (open it with its icon, or \ui{\textbackslash} for both rails) lists where the
candidate came from: the run and its recipe, the detection's score, other runs that found
the same place, the nearest family, artifact likelihood and the candidate's history. Where
nothing has computed a value the rail says ``not computed'' rather than showing a number.
The left-hand rail lists the queues and the candidates up next, with filters for channel,
score and status.

\screenshot{pages/review/review.inspector--4-evidence-rail (deep link)}{The evidence rail
open beside a candidate. Sections with no computed value say so in words.}

\subsection{Blind verification}\label{sec:review-blind}
In a blind queue the machine's opinion is hidden until you have given your verdict: the
score, family affinity and nearest families, the model's prediction, and the machine
sections of the evidence rail. The signal, the other channels and artifact likelihood are
never hidden. The reason is the same as for marking an exam without seeing the first
examiner's mark: if you can see what the machine said, your verdict stops being an
independent check of it. Window queues from training and from model verification are blind
by default, and take only \emph{interesting} and \emph{not interesting} plus a class. After
each verdict the \emph{Previous window, revealed} card shows your call beside the model's;
running agreement is shown only on \apppage{models.registry}, so it cannot steer the next
answer. In a queue the server marks blind, the score is not even sent to the browser.

\towrite{Screenshot of a blind window queue: the smoke walk's blind-verification states
(review.inspector--5-...) still point at fixture queue ids and show ``no review queue''.}

\subsection{The blind test queue}\label{sec:review-blind-test}
A queue built for a blind test of a trained model opens on its own page. It shows only the
window's raw trace in millivolts, shaded, with one window-length of context each side and its
length in words. There is no prediction, score, recording, channel, place in the recording
or queue rail, and a window shown twice (to measure how consistent you are) does not show
your first answer. The keys are \ui{I} interesting, \ui{N} not interesting, \ui{U}
\emph{can't tell} (stored as \emph{unsure}), \ui{A} artifact and \ui{Z} undo. Each answer is
written at once, and the queue resumes at the first unlabelled window next time.

\screenshot{pages/zzzzzzz_ah_blind_test/review.blind--fixup-ah-label-close-reopen-it-resumes-and-finishes}{The
blind test page: a window and nothing that could tip the answer.}

\subsection{Suspected artifacts}\label{sec:review-suspected}
A queue of suspected artifacts, sent from a Library family, adds a card that draws every
channel of the recording over the member's span on one shared millivolt axis. The channels
that flagged it carry their $r$, lag, amplitude ratio and chance test. The machine only
flags; you decide. This queue takes artifact, interesting, not interesting or unsure, but
not seed.

\pagefield{Where it leads}
\apppage{library.family} from a promotion or a nearest family; \apppage{discovery.seed} to
search for more of a new seed; \apppage{discovery.runs} from a run in the evidence rail;
\apppage{explore.cross-channel} from \ui{Open all channels in Explore}.

\begin{watchout}[Things the page does not keep]
The family you choose in the promotion panel is not written to the database: the entry joins
a family at the next regrouping in Library. A class is written only as the verdict it
implies; the class name itself is a label in this browser tab. Turning \ui{blind for this
queue} off or on is remembered in this browser only and is not stored with your verdicts,
although its popover says it is. \ui{Edit span in Explore} opens a page that still runs on
example data (\apppage{explore.span-edit}). Artifact likelihood reads ``not computed'' on every
queue except suspected artifacts.
\end{watchout}

\begin{underthehood}[Inspector]
Client: \code{webui/client/src/review/index.tsx} (routing), \code{Inspector.tsx},
\code{Shell.tsx} (toolbar and rails), \code{parts.tsx} (cards), \code{keys.ts} (keymap),
\code{store.ts} (the session undo stack, updated only after the server accepts a write),
\code{BlindTest.tsx}, \code{SuspectedArtifact.tsx}. Reads and writes go through
\code{webui/client/src/api/review.ts} to \code{/api/review/...} in
\code{webui/server/review.py}, which calls \code{Working.review}. The queue kinds, and the
unit and table each one writes, are fixed in \code{Working/review/queues.py}
(\code{_DEFAULTS}): detection queues write \code{adjudications}, span, sequence and member
queues write \code{annotations}, window queues write \code{window_verdicts}.
\code{Working/review/verdicts.py} (\code{write_verdict}, \code{write_batch},
\code{undo_last}) refuses any verdict that would cross that line with a
\code{PermissionError}, shown as a red ``Not written'' card. Promotion is
\code{Working/review/promotion.py} (\code{promote}, \code{unpromote}). The five words are
\code{VERDICTS} in \code{Working/database/schema.py}.
\end{underthehood}

\pagesection{review.cluster}{Cluster}{Seed-search matches with their member strip}{unchecked}
The cluster page judges a group of candidates that are thought to be the same kind of event
in one gesture. It shows one member in context, a strip of every member drawn on one shared
millivolt scale and sorted by distance to the group's most typical member, and a tick box
beside each.

\pagefield{When you'd come here}
When a queue groups its candidates, for example matches of one seed, and you want to accept
or reject many similar shapes at once rather than one by one.

\towrite{Screenshot: every cluster screenshot under webui/screenshots/pages/review/ predates the
live wiring and shows a skeleton or an error; none shows a working cluster.}

\pagefield{Doing things here}
\begin{enumerate}
  \item Look along the member strip. Members farther from the typical member than the
    cohesion limit are drawn in amber and left out of the batch by default. \ui{Include all}
    puts them back; \ui{Exclude flagged} takes them out again.
  \item Untick any member you do not want the verdict to cover. Click a member to show it in
    the context card.
  \item Give a verdict. With \ui{Verdict for the} \emph{N} \ui{included members} ticked, it
    covers every ticked member, up to 50; untick it to judge only the member shown.
  \item A seed here promotes one member, the one nearest the typical shape, and marks the
    rest interesting.
  \item \ui{Ctrl Z} undoes the whole batch in one step, because the server stores the batch
    as one act.
\end{enumerate}

\pagefield{Where it leads}
Back to the next unit of the queue in \apppage{review.inspector}.

\begin{watchout}[No live queue makes clusters yet]
The page reads its clusters from the bridge, but the bridge builds a cluster only from items
that carry a cluster number, and no queue source in \code{Working/review/} sets one. A real
queue therefore has no clusters, and asking for one gives a message saying so. The cluster
the bridge would describe also lacks the cohesion limit and family the page draws. Treat this
page as not yet usable.
\end{watchout}

\begin{underthehood}[Cluster]
Client: \code{webui/client/src/review/ClusterView.tsx}. Read: \code{GET
/api/review/queues/{q}/cluster/{n}}; writes: \code{/batch} for a chosen set of members,
or \code{/cluster/{n}/accept} and \code{/reject} for a whole untagged cluster (accept
writes \emph{interesting}, reject \emph{not interesting}). Clusters come from
\code{_clusters} in \code{webui/server/review.py}. The short link \code{#/review/cluster/<n>}
still resolves against fixture clusters (\code{clusterQueue} in \code{api/review.ts}).
\end{underthehood}
